\appendix

\section{Proof Details}
\label{sec:proof-details}

\subsection{Boundary cases and the architecture count}
\label{app:step-001}

\begin{proposition}[Exact empty-class closure]
\label{prop:step-001-empty}
Under Assumption~\ref{assump:source-regime} and the convention
$\operatorname{dc}(\varnothing)=0$, if $\mathcal H=\varnothing$, then
\[
\operatorname{dc}(\mathcal H)=0\leq 7TSd.
\]
\end{proposition}

\begin{proof}
The defining convention gives
$\operatorname{dc}(\mathcal H)=\operatorname{dc}(\varnothing)=0$.
Assumption~\ref{assump:source-regime} gives $T\geq 1$, $d\geq 0$, and a
nonempty sum $S=\sum_{i=1}^L n_i n_{i-1}$ of positive integer products.
Hence $7TSd\geq 0$, which proves the inequality.
\end{proof}

\begin{lemma}[Zero-dimensional tie-resolved closure]
\label{lem:step-001-zero}
Under Assumptions~\ref{assump:source-regime} and
\ref{assump:tie-resolved-confident-map}, and the conventions
$\mathbb R^0=\{0\}$, $\langle 0,0\rangle=0$, and
$\operatorname{sgn}_{\tau}(0)=\tau$, if $\mathcal H\neq\varnothing$ and
$d=0$, then every $h\in\mathcal H$ satisfies $h(x)=\tau$ for all
$x\in\mathcal X$, and
\[
\operatorname{dc}(\mathcal H)=0=7TSd.
\]
\end{lemma}

\begin{proof}
Since $\mathbb R^0$ is the singleton $\{0\}$, there is exactly one map
from $\mathcal X$ to $\mathbb R^0$, namely $\phi_0$ with
$\phi_0(x)=0$ for every $x$.  Therefore every probability law on such maps,
including the law $\mathcal P$ in Assumption~\ref{assump:tie-resolved-confident-map},
assigns probability one to $\phi_0$.  There is also exactly one candidate
separator $w\in\mathbb R^0$, namely $w=0$.

Fix any $h\in\mathcal H$.  Because $n\geq 1$, the domain
$\mathcal X=\{-1,+1\}^n$ is nonempty, so the confident-map premise may be
instantiated at any distribution on $\mathcal X$.  Its success event for this
fixed $h$ is
\[
E_h:=\left\{\phi:\ \exists w\in\mathbb R^0\ \forall x\in\mathcal X,
\operatorname{sgn}_{\tau}(\langle w,\phi(x)\rangle)=h(x)\right\}.
\]
For the sole possible map $\phi_0$ and the sole separator,
\[
\begin{aligned}
\phi_0\in E_h
&\iff \forall x\in\mathcal X,
\operatorname{sgn}_{\tau}(\langle 0,0\rangle)=h(x)\\
&\iff \forall x\in\mathcal X,\ h(x)=\tau.
\end{aligned}
\]
Thus $E_h$ is deterministic: its $\mathcal P$-probability is either zero
or one.  Assumption~\ref{assump:tie-resolved-confident-map} gives
$\mathcal P(E_h)\geq 1/2$, so $\mathcal P(E_h)=1$ and
$h(x)=\tau$ for every $x$.  Since $h$ was arbitrary, every member of
$\mathcal H$ is the constant $\tau$ classifier.

The same map $\phi_0$, together with the separator $u_h=0$ for every
$h\in\mathcal H$, therefore satisfies
\[
\operatorname{sgn}_{\tau}(\langle u_h,\phi_0(x)\rangle)
=\operatorname{sgn}_{\tau}(0)=\tau=h(x)
\qquad\text{for all }h\in\mathcal H,\ x\in\mathcal X.
\]
Dimension $q=0$ is consequently feasible in the definition of
$\operatorname{dc}(\mathcal H)$.  Since that definition minimizes over
$q\in\mathbb Z_{\geq 0}$, it follows that $\operatorname{dc}(\mathcal H)=0$.
Finally, $d=0$ gives $7TSd=0$, proving exact equality.
\end{proof}

\begin{proposition}[First-layer structural bound]
\label{prop:step-001-architecture}
Under Assumption~\ref{assump:source-regime}, if
$\mathcal H\neq\varnothing$ and $d\neq 0$, then
\[
d\geq 1,\qquad S\geq n\geq 1,\qquad T\geq 1,\qquad S\geq 1.
\]
Moreover, if $L=1$, then $S=n$.
\end{proposition}

\begin{proof}
Assumption~\ref{assump:source-regime} states that
$d\in\mathbb Z_{\geq 0}$.  Hence the branch condition $d\neq 0$ implies
$d\geq 1$.  The same assumption gives $n,L,T\in\mathbb Z_{\geq 1}$, so
$n\geq 1$ and $T\geq 1$.

For every allowed depth, $n_1\geq 1$: when $L\geq 2$, it is a positive
hidden width, while when $L=1$, it is the output width $n_1=n_L=1$.
The first summand in the definition of $S$ therefore obeys
\[
n_1n_0=n_1n\geq n.
\]
All summands in $S=\sum_{i=1}^L n_i n_{i-1}$ are positive, so
\[
S\geq n_1n_0\geq n\geq 1.
\]
This proves both $S\geq n$ and $S\geq 1$.  If $L=1$, the sum has only
its first term and $n_1=n_L=1$; hence
\[
S=n_1n_0=1\cdot n=n.
\]
Thus the same count covers the depth-one boundary without presupposing a
hidden layer.
\end{proof}

\begin{proof}[Boundary reduction]
The case split is exhaustive.  If $\mathcal H=\varnothing$,
Proposition~\ref{prop:step-001-empty} proves the target exactly.  Otherwise
$\mathcal H\neq\varnothing$.  In that branch, if $d=0$,
Lemma~\ref{lem:step-001-zero} proves
$\operatorname{dc}(\mathcal H)=0=7TSd$ with the fixed tie convention.  If
instead $d\neq 0$, Proposition~\ref{prop:step-001-architecture} gives
\[
\mathcal H\neq\varnothing,\qquad d\geq 1,\qquad S\geq n\geq 1,
\qquad T,S\geq 1,
\]
and explicitly includes the $L=1$ case.  These are exactly the boundary
conclusions and remaining-branch structural controls needed below.
\end{proof}

\subsection{A VC ceiling from the exact learner}
\label{app:step-002}

\begin{lemma}[Sampled-label measurability of exact SGD]
\label{lem:step-002-label-access}
Under Assumption~\ref{assump:source-regime}, suppose $Z\subseteq\mathcal X$
has $m=2T$ points and, for every $b\in\{-1,+1\}^{Z}$, a fixed target
$h_b\in\mathcal H$ satisfies $h_b(z)=b(z)$ for every $z\in Z$.  Let
\[
B=(B_z)_{z\in Z},\qquad B_z\stackrel{\mathrm{iid}}{\sim}
\operatorname{Unif}\{-1,+1\},
\]
let $X_0,\ldots,X_{T-1}$ be iid uniform on $Z$, let $X$ be an independent
uniform point of $Z$, and draw the source Gaussian initialization
$\Theta^{(0)}$ independently of $B,X_0,\ldots,X_{T-1},X$.  Run the exact
source recursion with target $h_B$.  Then, for every $0\leq t\leq T$,
\[
\Theta^{(t)}\text{ is measurable with respect to }
\mathscr F_t:=\sigma\!\left(\Theta^{(0)},(X_s,B_{X_s})_{0\leq s<t}\right),
\]
and the returned prediction at $X$,
\[
\widehat g_B(X):=\widehat h_{\mathcal D_Z,h_B}(X),
\qquad \mathcal D_Z:=\operatorname{Unif}(Z),
\]
is measurable with respect to
\[
\mathscr F:=\sigma\!\left(X,\Theta^{(0)},(X_s,B_{X_s})_{s=0}^{T-1}\right).
\]
\end{lemma}

\begin{proof}
For a parameter value $\theta$, an input $x\in\mathcal X$, and a label
$y\in\{-1,+1\}$, define the exact update map
\[
U(\theta,x,y):=\theta-\eta\nabla^{\mathrm{src}}_\theta
\ell\!\left(y f_\theta(x)\right).
\]
All conventions in this expression are fixed before the target is selected.
In particular, the fixed source choice at every ReLU kink makes $U$
single-valued; it cannot encode an unsampled value of the target.

Because every $X_t$ lies in $Z$, the target label passed to update $t$ is
exactly $h_B(X_t)=B_{X_t}$.  The recursion therefore has the pathwise form
\[
\Theta^{(t+1)}=U\!\left(\Theta^{(t)},X_t,B_{X_t}\right).
\]
At $t=0$, $\Theta^{(0)}$ is $\mathscr F_0$-measurable.  If
$\Theta^{(t)}$ is $\mathscr F_t$-measurable, the displayed update shows that
$\Theta^{(t+1)}$ is measurable with respect to
\[
\sigma(\mathscr F_t,X_t,B_{X_t})=\mathscr F_{t+1}.
\]
Induction proves the iterate claim for all $0\leq t\leq T$.  This induction
is pathwise and permits arbitrary repetitions among the $X_t$'s.

Consequently every term $f_{\Theta^{(t)}}(X)$ in
\[
\sum_{t=\lceil T/2\rceil}^{T}f_{\Theta^{(t)}}(X)
\]
is $\mathscr F$-measurable.  Applying the fixed map
$\operatorname{sgn}_{\tau}$ yields an $\mathscr F$-measurable value in
$\{-1,+1\}$, even when the aggregate is zero, because that case returns the
fixed label $\tau$.  Thus $\widehat g_B(X)$ is determined by exactly the
information listed in $\mathscr F$, and no unsampled coordinate of $B$ enters
the recursion or the output.
\end{proof}

\begin{lemma}[Fair unseen bit]
\label{lem:step-002-unseen-label}
Under Assumption~\ref{assump:source-regime} and the experiment and notation of
Lemma~\ref{lem:step-002-label-access}, define
\[
E:=\{X\notin\{X_0,\ldots,X_{T-1}\}\}.
\]
Then, for each $s\in\{-1,+1\}$,
\[
\Pr(B_X=s\mid\mathscr F)=\frac12
\quad\text{almost surely on }E,
\]
and
\[
\Pr\!\left(\widehat g_B(X)\ne B_X\mid\mathscr F\right)
=\frac12
\quad\text{almost surely on }E.
\]
\end{lemma}

\begin{proof}
Condition first on concrete values
\[
X=x,\qquad (X_0,\ldots,X_{T-1})=(x_0,\ldots,x_{T-1}).
\]
On $E$, $x$ differs from every $x_t$.  Let
$I=\{x_0,\ldots,x_{T-1}\}$; repetitions merely make $|I|<T$ and do not
add a new label coordinate.  The observed label vector
$(B_{x_t})_{t=0}^{T-1}$ is measurable with respect to $(B_z)_{z\in I}$.
Since the coordinates $(B_z)_{z\in Z}$ are mutually independent fair signs
and $x\notin I$, the coordinate $B_x$ remains a fair sign after conditioning
on every observed label.  The initialization $\Theta^{(0)}$ was drawn
independently of the complete label vector, so conditioning on it does not
change this conclusion.  Averaging this pointwise conditional statement over
the conditioned input values gives
\[
\Pr(B_X=s\mid\mathscr F)=\frac12
\quad\text{on }E.
\]
The event $E$ itself is $\mathscr F$-measurable because it depends only on
$X,X_0,\ldots,X_{T-1}$.

By Lemma~\ref{lem:step-002-label-access}, $\widehat g_B(X)$ is
$\mathscr F$-measurable and takes one fixed value in $\{-1,+1\}$ under this
conditioning.  Therefore, on $E$,
\[
\begin{aligned}
\Pr\!\left(\widehat g_B(X)\ne B_X\mid\mathscr F\right)
&=\sum_{s\in\{-1,+1\}}
\mathbf 1\{\widehat g_B(X)\ne s\}
\Pr(B_X=s\mid\mathscr F)\\
&=\frac12.
\end{aligned}
\]
If the aggregate score is zero, the first factor still contains the fixed
binary value $\tau$, so the same calculation applies without a no-tie or
margin assumption.
\end{proof}

\begin{lemma}[Finite-horizon avoidance]
\label{lem:step-002-avoidance}
Under Assumption~\ref{assump:source-regime}, if $X,X_0,\ldots,X_{T-1}$ are
independent and uniform on a set $Z$ of size $2T$, then
\[
\Pr\!\left[X\notin\{X_0,\ldots,X_{T-1}\}\right]
=\left(1-\frac1{2T}\right)^T\geq\frac12.
\]
For $T=1$, both the exact probability and the displayed lower bound equal
$1/2$.
\end{lemma}

\begin{proof}
Conditional on any fixed $X=x\in Z$, each $X_t$ avoids $x$ with probability
\[
1-\frac1{|Z|}=1-\frac1{2T}.
\]
The $T$ training draws are independent, including when their realized values
repeat, so multiplication gives
\[
\Pr(E\mid X=x)=\left(1-\frac1{2T}\right)^T.
\]
This expression does not depend on $x$, proving the equality after averaging
over $X$.

For completeness, for every integer $k\geq 1$ and $u\in[0,1]$,
\[
(1-u)^k\geq 1-ku.
\]
Indeed, equality holds for $k=1$; if it holds for $k$, multiplication by
$1-u\geq 0$ gives
\[
(1-u)^{k+1}\geq(1-ku)(1-u)
=1-(k+1)u+ku^2\geq1-(k+1)u.
\]
Applying this induction with $k=T$ and $u=1/(2T)\in(0,1]$ yields
\[
\left(1-\frac1{2T}\right)^T
\geq1-\frac{T}{2T}=\frac12.
\]
When $T=1$, the exact product is $1-1/2=1/2$, so the boundary case incurs
no slack.
\end{proof}

\begin{proposition}[Random-label average-risk lower bound]
\label{prop:step-002-average-risk}
Under Assumption~\ref{assump:source-regime} and Lemmas
\ref{lem:step-002-label-access}, \ref{lem:step-002-unseen-label}, and
\ref{lem:step-002-avoidance}, if $Z\subseteq\mathcal X$ has $2T$ points and
every labeling $b\in\{-1,+1\}^{Z}$ has a fixed representative
$h_b\in\mathcal H$ with $h_b|_Z=b$, then, for
$\mathcal D_Z=\operatorname{Unif}(Z)$,
\[
2^{-2T}\sum_{b\in\{-1,+1\}^{Z}}
\mathbb E_{\Theta^{(0)},X_{0:T-1}\stackrel{\mathrm{iid}}{\sim}\mathcal D_Z}
\!\left[\mathcal L_{\mathcal D_Z,h_b}
(\widehat h_{\mathcal D_Z,h_b})\right]\geq\frac14.
\]
\end{proposition}

\begin{proof}
In the joint experiment of Lemma~\ref{lem:step-002-label-access}, condition
on $B=b$, the initialization, and the ordered training sample.  The
independent point $X\sim\mathcal D_Z$ converts the population risk into an
indicator expectation.  Since $X\in Z$ and $h_b|_Z=b$,
\[
\mathcal L_{\mathcal D_Z,h_b}(\widehat h_{\mathcal D_Z,h_b})
=\mathbb E_X\!\left[
\mathbf 1\{\widehat h_{\mathcal D_Z,h_b}(X)\ne b(X)\}
\,\middle|\,b,\Theta^{(0)},X_{0:T-1}\right].
\]
Taking expectation over the uniform labeling $B$, initialization, and
training sample, and using that $B$ has exactly $2^{2T}$ equiprobable values,
gives the exact identity
\[
\begin{aligned}
&2^{-2T}\sum_{b\in\{-1,+1\}^{Z}}
\mathbb E_{\Theta^{(0)},X_{0:T-1}}
\!\left[\mathcal L_{\mathcal D_Z,h_b}
(\widehat h_{\mathcal D_Z,h_b})\right]\\
&\hspace{25mm}=\Pr\!\left[\widehat g_B(X)\ne B_X\right].
\end{aligned}
\]
Because $E\in\mathscr F$, the tower property and
Lemma~\ref{lem:step-002-unseen-label} yield
\[
\begin{aligned}
\Pr\!\left[\widehat g_B(X)\ne B_X\right]
&\geq\mathbb E\!\left[\mathbf 1_E\mathbf 1\{\widehat g_B(X)\ne B_X\}\right]\\
&=\mathbb E\!\left[\mathbf 1_E
\Pr\!\left(\widehat g_B(X)\ne B_X\mid\mathscr F\right)\right]\\
&=\frac12\Pr(E).
\end{aligned}
\]
Lemma~\ref{lem:step-002-avoidance} gives $\Pr(E)\geq 1/2$, and hence the
finite average is at least $1/4$.  At $T=1$, $\Pr(E)=1/2$ exactly, so this
argument still gives the exact lower bound $1/4$.
\end{proof}

\begin{proposition}[VC ceiling from universal exact-SGD success]
\label{prop:step-002-vc}
Under Assumptions~\ref{assump:source-regime} and
\ref{assump:universal-sgd-success}, and on the remaining branch identified by
Proposition~\ref{prop:step-001-architecture},
\[
\operatorname{VC}(\mathcal H)<2T.
\]
\end{proposition}

\begin{proof}
Suppose instead that $\operatorname{VC}(\mathcal H)\geq 2T$.  Then there is
a set $Z\subseteq\mathcal X$ of exactly $2T$ points shattered by $\mathcal H$.
For each of the finitely many $b\in\{-1,+1\}^{Z}$, shattering supplies at
least one $h\in\mathcal H$ with $h|_Z=b$.  Fix one such representative and
call it $h_b$.

For each fixed labeling $b$, define
\[
R_b:=\mathbb E_{\Theta^{(0)},X_{0:T-1}\stackrel{\mathrm{iid}}{\sim}\mathcal D_Z}
\!\left[\mathcal L_{\mathcal D_Z,h_b}
(\widehat h_{\mathcal D_Z,h_b})\right],
\qquad \mathcal D_Z=\operatorname{Unif}(Z).
\]
Proposition~\ref{prop:step-002-average-risk} proves
\[
2^{-2T}\sum_{b\in\{-1,+1\}^{Z}}R_b\geq\frac14.
\]
This is an average over a finite set, so at least one fixed labeling $b_\star$
satisfies $R_{b_\star}\geq1/4$.  The auxiliary random labeling used to
establish the average is now gone: $\mathcal D_Z$ and
$h_{b_\star}\in\mathcal H$ are fixed, and $R_{b_\star}$ averages only over
the exact Gaussian initialization and the $T$ iid training draws specified
in the setting.

The architecture, step size $\eta$, and horizon $T$ were fixed before this
witness distribution and target were chosen.  Therefore
Assumption~\ref{assump:universal-sgd-success} applies directly to this same
fixed pair $(\mathcal D_Z,h_{b_\star})$ and the same exact learner, giving
\[
R_{b_\star}\leq\varepsilon<\frac14.
\]
This contradicts $R_{b_\star}\geq 1/4$.  Thus no $2T$-point subset of
$\mathcal X$ is shattered, and $\operatorname{VC}(\mathcal H)<2T$.
When $T=1$, the lower bound remains $1/4$ and the strict assumption
$\varepsilon<1/4$ gives the same contradiction.
\end{proof}

\begin{proof}[VC-dimension conclusion]
Proposition~\ref{prop:step-001-architecture} places the proof on the
nondegenerate branch.  Lemma~\ref{lem:step-002-label-access} shows from the
exact update and fixed kink convention that the learner uses only sampled
target labels.  Lemma~\ref{lem:step-002-unseen-label} gives conditional error
$1/2$ on an unsampled test point, including repeated inputs and zero aggregate
scores.  Lemma~\ref{lem:step-002-avoidance} gives the exact probability
$(1-1/(2T))^T\geq 1/2$, and Proposition~\ref{prop:step-002-average-risk}
converts these facts to the finite average risk lower bound.  The finite
selection and the strict comparison with $\varepsilon<1/4$ in
Proposition~\ref{prop:step-002-vc} then yield the displayed VC ceiling.
\end{proof}

\subsection{Finite-class counting and the repetition budget}
\label{app:step-003}

\begin{lemma}[Finite-domain growth bound]
\label{lem:step-003-growth}
Let $A$ be a finite set of cardinality $m$, and let
$\mathcal G\subseteq\{-1,+1\}^{A}$ be nonempty.  If
$q=\operatorname{VC}(\mathcal G)$, then $0\leq q\leq m$ and
\[
|\mathcal G|\leq\sum_{j=0}^{q}\binom mj.
\]
If $q=0$, then in fact $|\mathcal G|=1$.
\end{lemma}

\begin{proof}
A shattered set is a subset of $A$, so $0\leq q\leq m$.  We prove the
cardinality inequality in the slightly more general form that every nonempty
$\mathcal G\subseteq\{-1,+1\}^{A}$ satisfying
$\operatorname{VC}(\mathcal G)\leq q$, for an integer $0\leq q\leq m$,
obeys
\begin{equation}
|\mathcal G|\leq\sum_{j=0}^{q}\binom mj.
\label{eq:step-003-growth-induction}
\end{equation}
Use induction on $m$, with the convention $\binom ab=0$ when $b<0$ or
$b>a$.

If $m=0$, there is exactly one function on $A=\varnothing$, so nonemptiness
gives $|\mathcal G|=1=\binom00$.  For arbitrary $m$, if $q=0$, two
distinct functions in $\mathcal G$ would differ at some $a\in A$.  Those two
functions would realize both labels on $\{a\}$, thereby shattering that
singleton and contradicting $\operatorname{VC}(\mathcal G)\leq0$.  Hence
$|\mathcal G|\leq1=\binom m0$.  When $\mathcal G$ is nonempty and its VC
dimension is exactly zero, this also proves $|\mathcal G|=1$.  If $q=m$,
then the trivial count of all binary functions gives
\[
|\mathcal G|\leq2^m=\sum_{j=0}^{m}\binom mj.
\]

It remains to prove \eqref{eq:step-003-growth-induction} when $m\geq1$ and
$1\leq q<m$.  Fix $a\in A$, put $B=A\setminus\{a\}$, and, for
$s\in\{-1,+1\}$, define the restriction class
\[
\mathcal G_s:=\{g|_B:g\in\mathcal G,\ g(a)=s\}
\subseteq\{-1,+1\}^{B}.
\]
Let
\[
\mathcal U:=\mathcal G_{+1}\cup\mathcal G_{-1},
\qquad
\mathcal P_0:=\mathcal G_{+1}\cap\mathcal G_{-1}.
\]
A function in $\mathcal G$ is uniquely determined by its restriction to $B$
and its value at $a$.  Inclusion--exclusion therefore gives the exact
decomposition
\begin{equation}
|\mathcal G|
=|\mathcal G_{+1}|+|\mathcal G_{-1}|
=|\mathcal U|+|\mathcal P_0|.
\label{eq:step-003-decomposition}
\end{equation}

If $\mathcal U$ shatters $D\subseteq B$, then each labeling of $D$ has an
extension in $\mathcal G$, so $\mathcal G$ also shatters $D$.  Thus
$\operatorname{VC}(\mathcal U)\leq q$.  If $\mathcal P_0$ is nonempty and
shatters $D\subseteq B$, then for each labeling of $D$ its realizing
restriction in $\mathcal P_0$ has one extension in $\mathcal G$ with label
$+1$ at $a$ and another with label $-1$ there.  Hence $\mathcal G$ shatters
$D\cup\{a\}$, which shows $\operatorname{VC}(\mathcal P_0)\leq q-1$.  If
$\mathcal P_0$ is empty, the bound below is immediate for it.

Both $\mathcal U$ and $\mathcal P_0$ are classes on the $(m-1)$-point set
$B$.  Applying the induction hypothesis, and using the zero bound when
$\mathcal P_0=\varnothing$, gives
\[
|\mathcal U|\leq\sum_{j=0}^{q}\binom{m-1}{j},
\qquad
|\mathcal P_0|\leq\sum_{j=0}^{q-1}\binom{m-1}{j}.
\]
For completeness, partitioning the $j$-element subsets of an $m$-element set
according to whether they contain one fixed element gives
\[
\binom mj=\binom{m-1}{j}+\binom{m-1}{j-1}.
\]
Combining this identity with \eqref{eq:step-003-decomposition} yields
\[
\begin{aligned}
|\mathcal G|
&\leq\sum_{j=0}^{q}\binom{m-1}{j}
+\sum_{j=0}^{q-1}\binom{m-1}{j}\\
&=\sum_{j=0}^{q}
\left(\binom{m-1}{j}+\binom{m-1}{j-1}\right)
=\sum_{j=0}^{q}\binom mj.
\end{aligned}
\]
This closes the induction and proves the lemma.
\end{proof}

\begin{lemma}[Binomial-sum and numerical logarithm estimates]
\label{lem:step-003-binomial}
For integers $m,q$ satisfying $1\leq q\leq m$,
\[
\sum_{j=0}^{q}\binom mj\leq\left(\frac{em}{q}\right)^q.
\]
Moreover,
\[
\log_2 e<\frac32.
\]
\end{lemma}

\begin{proof}
Set $a=q/m\in(0,1]$.  For every $0\leq j\leq q$, one has
$a^j\geq a^q$.  Expanding the product of $m$ copies of $1+a$, and grouping
terms by the number $j$ of copies from which $a$ is selected, gives
\[
(1+a)^m=\sum_{j=0}^{m}\binom mj a^j.
\]
All terms are nonnegative, and hence
\begin{equation}
\begin{aligned}
\sum_{j=0}^{q}\binom mj
&\leq a^{-q}\sum_{j=0}^{q}\binom mj a^j\\
&\leq a^{-q}(1+a)^m.
\end{aligned}
\label{eq:step-003-binomial-weight}
\end{equation}
For $a\geq0$, the exponential series gives
\[
e^a=\sum_{k=0}^{\infty}\frac{a^k}{k!}\geq1+a.
\]
Consequently $(1+a)^m\leq e^{am}=e^q$.  Substituting this in
\eqref{eq:step-003-binomial-weight} and using
$a^{-q}=(m/q)^q$ proves
\[
\sum_{j=0}^{q}\binom mj\leq\left(\frac{em}{q}\right)^q.
\]

It remains to verify the numerical logarithm bound without a decimal
approximation.  For every $k\geq4$,
\[
k!\geq24\,4^{k-4},
\]
because $4!=24$ and each subsequent factor is at least $4$.  Hence
\[
\begin{aligned}
e=\sum_{k=0}^{\infty}\frac1{k!}
&\leq1+1+\frac12+\frac16
+\frac1{24}\sum_{j=0}^{\infty}4^{-j}\\
&=1+1+\frac12+\frac16+\frac1{18}
=\frac{49}{18}.
\end{aligned}
\]
Since
\[
\left(\frac{49}{18}\right)^2
=\frac{2401}{324}<\frac{2592}{324}=8,
\]
we have $e\leq49/18<\sqrt8=2^{3/2}$.  Monotonicity of the base-two
logarithm now gives $\log_2e<3/2$.
\end{proof}

\begin{proposition}[Explicit repetition budget]
\label{prop:step-003-budget}
Under Assumption~\ref{assump:source-regime},
Propositions~\ref{prop:step-001-architecture} and
\ref{prop:step-002-vc}, and Lemmas~\ref{lem:step-003-growth} and
\ref{lem:step-003-binomial}, on the nondegenerate branch define
\[
N:=|\mathcal X|=2^n,
\qquad M:=|\mathcal H|,
\qquad v:=\operatorname{VC}(\mathcal H),
\qquad r:=\left\lceil\log_2(2M)\right\rceil.
\]
Then $r\leq7TS$.
\end{proposition}

\begin{proof}
Proposition~\ref{prop:step-001-architecture} gives
$\mathcal H\neq\varnothing$, so $M\geq1$ and $r$ is well-defined.  It also
gives
\begin{equation}
n\geq1,\qquad S\geq n,\qquad T,S\geq1.
\label{eq:step-003-structural-inputs}
\end{equation}
Proposition~\ref{prop:step-002-vc} gives
\begin{equation}
v<2T.
\label{eq:step-003-vc-input}
\end{equation}

First suppose $v=0$.  Applying Lemma~\ref{lem:step-003-growth} to
$A=\mathcal X$ and $\mathcal G=\mathcal H$ gives $M=1$.  Therefore
\[
r=\left\lceil\log_2 2\right\rceil=1\leq7TS,
\]
where the last inequality follows from $T,S\geq1$.  This handles the
singular case in which $(eN/v)^v$ is unavailable.

Now suppose $v\geq1$.  Lemma~\ref{lem:step-003-growth} gives $v\leq N$ and
\[
M\leq\sum_{j=0}^{v}\binom Nj.
\]
Lemma~\ref{lem:step-003-binomial} therefore applies, including when $v=N$,
and yields
\[
M\leq\left(\frac{eN}{v}\right)^v.
\]
Taking base-two logarithms, using $N=2^n$ and $\log_2v\geq0$, and then
using \eqref{eq:step-003-vc-input} and $\log_2e<3/2$, gives
\begin{equation}
\begin{aligned}
\log_2M
&\leq v\log_2\!\left(\frac{eN}{v}\right)\\
&=v\left(n+\log_2e-\log_2v\right)\\
&\leq v\left(n+\log_2e\right)\\
&\leq2T\left(n+\frac32\right)\\
&\leq5Tn.
\end{aligned}
\label{eq:step-003-log-cardinality}
\end{equation}
The last inequality is explicit: $n\geq1$ implies
\[
n+\frac32\leq\frac52n.
\]

For every real $y$, $\lceil y\rceil\leq y+1$.  Hence
\eqref{eq:step-003-log-cardinality} accounts for the ceiling as
\begin{equation}
\begin{aligned}
r
&=\left\lceil1+\log_2M\right\rceil\\
&\leq\log_2M+2\\
&\leq5Tn+2\\
&\leq5TS+2TS\\
&=7TS.
\end{aligned}
\label{eq:step-003-final-budget}
\end{equation}
Here $5Tn\leq5TS$ follows from $n\leq S$, while $2\leq2TS$ follows from
$T,S\geq1$.  Thus neither $n$ nor the additive ceiling contribution remains
in the exported bound.

The corner $T=1,n=1$ is included explicitly: then $N=2$ and
\eqref{eq:step-003-vc-input} forces $v\in\{0,1\}$, so one of the two branches
above applies.  In the $v\geq1$ branch,
$n+3/2=(5/2)n$; if also $S=1$, then $2=2TS$.  Thus the two potentially
tight elementary absorptions in \eqref{eq:step-003-log-cardinality}--
\eqref{eq:step-003-final-budget} remain valid at the smallest allowed values.
\end{proof}

\begin{proof}[Derivation of the repetition budget]
Proposition~\ref{prop:step-001-architecture} supplies nonemptiness,
$S\geq n\geq1$, and $T,S\geq1$, while
Proposition~\ref{prop:step-002-vc} supplies $v<2T$.
Lemma~\ref{lem:step-003-growth} proves both the finite-domain growth bound and
the exact $v=0\Rightarrow M=1$ boundary.  Lemma~\ref{lem:step-003-binomial}
proves the binomial estimate and $\log_2e<3/2$.
Proposition~\ref{prop:step-003-budget} then pays the ceiling by
$2\leq2TS$ and eliminates $n$ using $n\leq S$, yielding
\[
r=\left\lceil\log_2(2|\mathcal H|)\right\rceil\leq7TS.
\]
\end{proof}

\subsection{From target-wise confidence to one common tuple}
\label{app:step-004}

\begin{lemma}[Independent amplification for one target]
\label{lem:step-004-single-target}
Under Assumption~\ref{assump:tie-resolved-confident-map} and
Proposition~\ref{prop:step-003-budget}, fix the single law $\mathcal P$
supplied before all distributions and targets, and let
\[
r=\left\lceil\log_2(2|\mathcal H|)\right\rceil\geq1.
\]
If $\phi_1,\ldots,\phi_r$ are iid with law $\mathcal P$, then, for every
fixed $h\in\mathcal H$,
\[
\Pr_{(\phi_1,\ldots,\phi_r)\sim\mathcal P^r}
\!\left[
\begin{array}{c}
\forall i\in\{1,\ldots,r\},\ \nexists w\in\mathbb R^d\ \text{such that}\\
\forall x\in\mathcal X,\quad
\operatorname{sgn}_{\tau}(\langle w,\phi_i(x)\rangle)=h(x)
\end{array}
\right]
\leq2^{-r}.
\]
\end{lemma}

\begin{proof}
Choose the one $\mathcal P$ occurring before the universal $\mathcal D$- and
$h$-quantifiers in Assumption~\ref{assump:tie-resolved-confident-map}.  For a
fixed $h\in\mathcal H$, define the exact one-map event
\begin{equation}
E_h:=\left\{\phi:\mathcal X\to\mathbb R^d:\
\exists w\in\mathbb R^d\ \forall x\in\mathcal X,\quad
\operatorname{sgn}_{\tau}(\langle w,\phi(x)\rangle)=h(x)\right\}.
\label{eq:step-004-success-event}
\end{equation}
The probability in the assumption makes $E_h$ measurable.  The event contains
no $\mathcal D$.  The finite hypercube $\mathcal X$ is nonempty, so
$\Delta(\mathcal X)$ contains a point mass and the universal distribution
quantifier is not vacuous.  Consequently, after the single law $\mathcal P$
has been fixed, that quantifier does not change the event or permit a new law,
and the premise gives
\begin{equation}
\mathcal P(E_h)\geq\frac12,
\qquad
\mathcal P(E_h^c)\leq\frac12
\label{eq:step-004-one-map-probability}
\end{equation}
for every fixed $h\in\mathcal H$.

For coordinate $i$, the event that block $i$ fails to represent $h$ is
$\{\phi_i\in E_h^c\}$.  Since all coordinates have the same law $\mathcal P$
and are independent, the defining rectangle identity for the product law and
\eqref{eq:step-004-one-map-probability} give
\begin{equation}
\begin{aligned}
\Pr_{\mathcal P^r}\!\left[\bigcap_{i=1}^{r}
\{\phi_i\in E_h^c\}\right]
&=\prod_{i=1}^{r}\mathcal P(E_h^c)\\
&=\mathcal P(E_h^c)^r\\
&\leq\left(\frac12\right)^r=2^{-r}.
\end{aligned}
\label{eq:step-004-product-failure}
\end{equation}
Membership in $E_h$ is exactly the full-domain event in the statement, so
\eqref{eq:step-004-product-failure} is the claimed failure bound.  In
particular, the calculation neither requires nor permits a law depending on
$h$, and it never replaces $\operatorname{sgn}_{\tau}$ by a strict-margin or
approximate event.
\end{proof}

\begin{proposition}[Common deterministic exact-cover tuple]
\label{prop:step-004-covering}
Under Assumption~\ref{assump:tie-resolved-confident-map},
Proposition~\ref{prop:step-003-budget}, and
Lemma~\ref{lem:step-004-single-target}, let
\[
M:=|\mathcal H|,
\qquad
r:=\left\lceil\log_2(2M)\right\rceil.
\]
Then $1\leq r\leq7TS$, and there are deterministic maps
$\phi_i^*:\mathcal X\to\mathbb R^d$, $1\leq i\leq r$, such that
\begin{equation}
\forall h\in\mathcal H\ \exists i\in\{1,\ldots,r\}\
\exists w\in\mathbb R^d\ \forall x\in\mathcal X,
\qquad
\operatorname{sgn}_{\tau}(\langle w,\phi_i^*(x)\rangle)=h(x).
\label{eq:step-004-common-cover}
\end{equation}
\end{proposition}

\begin{proof}
Proposition~\ref{prop:step-003-budget} gives $M\geq1$, the displayed
definition of the positive integer $r$, and $r\leq7TS$.  Draw one random
tuple $(\phi_1,\ldots,\phi_r)\sim\mathcal P^r$ from the product of the
single law fixed by Assumption~\ref{assump:tie-resolved-confident-map}.  For
each $h\in\mathcal H$, let
\begin{equation}
F_h:=\bigcap_{i=1}^{r}\{\phi_i\in E_h^c\},
\label{eq:step-004-uncovered-event}
\end{equation}
where $E_h$ is the exact event \eqref{eq:step-004-success-event}.  Thus $F_h$
is precisely the event that this tuple leaves $h$ uncovered.  By
Lemma~\ref{lem:step-004-single-target},
\begin{equation}
\Pr_{\mathcal P^r}(F_h)\leq2^{-r}
\qquad\text{for every }h\in\mathcal H.
\label{eq:step-004-target-failure}
\end{equation}

The events $F_h$ may be dependent because they are evaluated on the same
tuple; no independence across targets is needed.  Since $\mathcal H$ is
finite, the pointwise indicator inequality
\[
\mathbf 1_{\cup_{h\in\mathcal H}F_h}
\leq\sum_{h\in\mathcal H}\mathbf 1_{F_h}
\]
and \eqref{eq:step-004-target-failure} imply
\begin{equation}
\Pr_{\mathcal P^r}\!\left[\bigcup_{h\in\mathcal H}F_h\right]
\leq\sum_{h\in\mathcal H}\Pr_{\mathcal P^r}(F_h)
\leq M2^{-r}.
\label{eq:step-004-union-bound}
\end{equation}
The definition of $r$ gives
\begin{equation}
r\geq\log_2(2M),
\qquad
M2^{-r}\leq M2^{-\log_2(2M)}=\frac12<1.
\label{eq:step-004-ceiling-payment}
\end{equation}
Combining \eqref{eq:step-004-union-bound}--
\eqref{eq:step-004-ceiling-payment}, the complementary simultaneous-coverage
event
\[
G:=\bigcap_{h\in\mathcal H}F_h^c
\]
satisfies
\begin{equation}
\Pr_{\mathcal P^r}(G)
=1-\Pr_{\mathcal P^r}\!\left[\bigcup_{h\in\mathcal H}F_h\right]
\geq\frac12>0.
\label{eq:step-004-positive-probability}
\end{equation}
The empty event has probability zero, so
\eqref{eq:step-004-positive-probability} proves that $G$ contains at least one
tuple.  Fix one such realization and call it
$(\phi_1^*,\ldots,\phi_r^*)$.  This tuple is fixed once, before any target is
selected from $\mathcal H$.  For every $h\in\mathcal H$, membership in
$F_h^c$ gives some $i\in\{1,\ldots,r\}$ with $\phi_i^*\in E_h$; by the
definition \eqref{eq:step-004-success-event} of $E_h$, that membership gives
some $w\in\mathbb R^d$ satisfying \eqref{eq:step-004-common-cover}.

Thus the final quantifier order is
\[
\exists(\phi_1^*,\ldots,\phi_r^*)\ \forall h\in\mathcal H\
\exists i\in\{1,\ldots,r\}\ \exists w\in\mathbb R^d\
\forall x\in\mathcal X,
\]
not a separate tuple or law chosen after $h$.  No distribution appears in the
fixed tuple or in \eqref{eq:step-004-common-cover}, because the exact event
\eqref{eq:step-004-success-event} is independent of $\mathcal D$.
\end{proof}

\begin{proof}[Derivation of simultaneous deterministic coverage]
Proposition~\ref{prop:step-003-budget} supplies the positive integer
\[
r=\left\lceil\log_2(2|\mathcal H|)\right\rceil\leq7TS.
\]
Assumption~\ref{assump:tie-resolved-confident-map} supplies one law before all
targets, and Lemma~\ref{lem:step-004-single-target} proves that $r$ iid draws
from this law miss any fixed target with probability at most $2^{-r}$.
Proposition~\ref{prop:step-004-covering} pays the complete finite union through
$|\mathcal H|2^{-r}\leq1/2<1$ and fixes a tuple in the positive-probability
simultaneous-coverage event.  The resulting tuple is deterministic, common to
all targets, full-domain, and tie-resolved exactly as in the premise.
When $|\mathcal H|=1$, the definition gives $r=1$, and the construction
reduces to the primitive success event for the sole target, whose probability
is at least $1/2$.  Thus success exactly at the stated threshold is sufficient.
Any successful zero score remains valid because the same
$\operatorname{sgn}_{\tau}$ is used before and after fixing the tuple.
\end{proof}

\subsection{Direct-sum representation and dimension closure}
\label{app:step-005}

\begin{proposition}[Exact boundary routing and nondegenerate controls]
\label{prop:step-005-routing}
Under Assumption~\ref{assump:source-regime},
Propositions~\ref{prop:step-001-empty} and
\ref{prop:step-001-architecture}, and Lemma~\ref{lem:step-001-zero}, the
following exhaustive alternatives hold:
\begin{enumerate}
\item If $\mathcal H=\varnothing$, then
$\operatorname{dc}(\mathcal H)=0\leq7TSd$.
\item If $\mathcal H\neq\varnothing$ and $d=0$, then
$\operatorname{dc}(\mathcal H)=0=7TSd$.
\item If $\mathcal H\neq\varnothing$ and $d\neq0$, then
$d\geq1$, $S\geq n\geq1$, and $T,S\geq1$.
\end{enumerate}
\end{proposition}

\begin{proof}
The three cases are exhaustive because $\mathcal H$ is either empty or
nonempty and $d\in\mathbb Z_{\geq0}$.  In the first case,
Proposition~\ref{prop:step-001-empty} gives the exact dimension-zero
conclusion.  In the second case, Lemma~\ref{lem:step-001-zero} gives the
exact common zero-map conclusion with the fixed convention
$\operatorname{sgn}_{\tau}(0)=\tau$.  In either case the logarithmic quantity
$r$ and a block tuple are unnecessary, so no undefined expression is
introduced when $\mathcal H=\varnothing$.  In the third case,
Proposition~\ref{prop:step-001-architecture} supplies all displayed
positivity and structural controls.
\end{proof}

\begin{lemma}[Exact direct-sum score preservation]
\label{lem:step-005-score}
Under Assumption~\ref{assump:source-regime},
Proposition~\ref{prop:step-001-architecture}, and
Proposition~\ref{prop:step-004-covering}, let
$r=\lceil\log_2(2|\mathcal H|)\rceil$ and let
$(\phi_1^*,\ldots,\phi_r^*)$ be the deterministic tuple supplied by
Proposition~\ref{prop:step-004-covering}.  If
$\mathcal H\neq\varnothing$ and $d\neq0$, then there are a single map
$\Phi:\mathcal X\to\mathbb R^{rd}$ and vectors
$u_h\in\mathbb R^{rd}$, one for each $h\in\mathcal H$, such that for every
$h$ and $x$,
\[
\langle u_h,\Phi(x)\rangle
=\langle w_h,\phi_{i(h)}^*(x)\rangle,
\qquad
\operatorname{sgn}_{\tau}(\langle u_h,\Phi(x)\rangle)=h(x),
\]
where $i(h)$ and $w_h$ are one successful block and separator for $h$.
The map $\Phi$ is independent of $h$.
\end{lemma}

\begin{proof}
By Proposition~\ref{prop:step-001-architecture}, the branch has $d\geq1$.
Proposition~\ref{prop:step-004-covering} supplies a fixed tuple and, for each
$h$, a nonempty finite set
\[
I_h:=\left\{i\in\{1,\ldots,r\}:\exists w\in\mathbb R^d\ \forall x\in\mathcal X,
\ \operatorname{sgn}_{\tau}(\langle w,\phi_i^*(x)\rangle)=h(x)\right\}.
\]
Choose $i(h):=\min I_h$, and choose one witnessing vector
$w_h\in\mathbb R^d$ for this index.  This is a finite target-by-target
selection of witnesses; it changes neither the already fixed tuple nor the
feature map.

Define the concatenated map by
\[
\Phi(x):=(\phi_1^*(x),\ldots,\phi_r^*(x))
\in(\mathbb R^d)^r\cong\mathbb R^{rd}.
\]
This definition uses only the deterministic tuple and therefore is one common
target-independent map.  For each target define the zero-padded vector
\[
u_h:=(u_{h,1},\ldots,u_{h,r})\in(\mathbb R^d)^r,
\qquad
u_{h,j}:=
\begin{cases}
w_h,&j=i(h),\\
0_d,&j\neq i(h),
\end{cases}
\]
where $0_d$ is the zero vector in $\mathbb R^d$.  The Euclidean product inner
product gives, for every $x\in\mathcal X$,
\[
\begin{aligned}
\langle u_h,\Phi(x)\rangle
&=\sum_{j=1}^{r}\langle u_{h,j},\phi_j^*(x)\rangle\\
&=\langle w_h,\phi_{i(h)}^*(x)\rangle,
\end{aligned}
\]
because every term with $j\neq i(h)$ is
$\langle0_d,\phi_j^*(x)\rangle=0$.  Thus the score residual is explicitly
\[
\langle u_h,\Phi(x)\rangle
-\langle w_h,\phi_{i(h)}^*(x)\rangle=0.
\]
Applying the same fixed $\operatorname{sgn}_{\tau}$ to equal scores preserves
the successful block's exact label, including when that score is zero.  Hence
$\operatorname{sgn}_{\tau}(\langle u_h,\Phi(x)\rangle)=h(x)$ for every
$h\in\mathcal H$ and $x\in\mathcal X$.
\end{proof}

\begin{proposition}[Common exact representation and dimension bound]
\label{prop:step-005-dimension}
Under Assumption~\ref{assump:source-regime},
Propositions~\ref{prop:step-001-architecture},
\ref{prop:step-003-budget}, and \ref{prop:step-004-covering}, and
Lemma~\ref{lem:step-005-score}, if $\mathcal H\neq\varnothing$ and
$d\neq0$, then
\[
\operatorname{dc}(\mathcal H)\leq rd\leq7TSd,
\qquad
r=\left\lceil\log_2(2|\mathcal H|)\right\rceil.
\]
\end{proposition}

\begin{proof}
Lemma~\ref{lem:step-005-score} supplies one map
$\Phi:\mathcal X\to\mathbb R^{rd}$ and, for every target, a vector
$u_h\in\mathbb R^{rd}$ with exact tie-resolved signs on all of $\mathcal X$.
By the definition of $\operatorname{dc}$, this is a feasible representation
of dimension $rd$, so $\operatorname{dc}(\mathcal H)\leq rd$.
Proposition~\ref{prop:step-003-budget} gives $r\leq7TS$.  Since $d\geq1$
by Proposition~\ref{prop:step-001-architecture},
\[
7TSd-rd=(7TS-r)d\geq0,
\]
and therefore $rd\leq7TSd$.  The integer product $rd$ is a valid dimension
because $r,d\in\mathbb Z_{\geq0}$.
\end{proof}

\begin{corollary}[Conditional polynomial rate specialization]
\label{cor:step-005-polynomial-bridge}
Under Assumption~\ref{assump:source-regime},
Proposition~\ref{prop:step-005-routing}, and
Proposition~\ref{prop:step-005-dimension} on the nondegenerate branch, suppose
additionally that a separate result supplies an explicit polynomial expression
$p(S,T)$ satisfying $d\leq p(S,T)$ at the current setup, with all coefficients
and auxiliary dependence exposed and with no hidden dependence on $n$ or
$\eta$.  Then every branch satisfies
\[
\operatorname{dc}(\mathcal H)\leq7TSp(S,T),
\]
which is polynomial in $S,T$ whenever $p$ is polynomial in those variables.
\end{corollary}

\begin{proof}
The additional inequality is a local conditional hypothesis for this
corollary only; it is not used to establish
Proposition~\ref{prop:step-005-dimension}.  From
Assumption~\ref{assump:source-regime}, $T\geq1$, and the positive
architecture widths imply $S\geq1$; hence $7TS\geq0$.  The supplied
inequality and $d\geq0$ imply $p(S,T)\geq d\geq0$ at the current parameters.
Propositions~\ref{prop:step-005-routing} and
\ref{prop:step-005-dimension}, applied in the nondegenerate alternative,
jointly give the all-branch main inequality
\[
\operatorname{dc}(\mathcal H)\leq7TSd.
\]
Indeed, the first proposition closes $\mathcal H=\varnothing$ and $d=0$,
and the second closes the remaining branch.  The exact absorption relation is
\[
7TSp(S,T)-7TSd=7TS\bigl(p(S,T)-d\bigr)\geq0.
\]
Combining the all-branch main inequality with this relation gives
\[
\operatorname{dc}(\mathcal H)\leq7TSd\leq7TSp(S,T).
\]
No representation score, tie, probability, or horizon term is changed: on the
nondegenerate branch the same deterministic map from
Lemma~\ref{lem:step-005-score} is used, and on the null branches the exact maps
remain in force.  Only a scalar upper bound is substituted.  If the separate
inequality is unavailable, this corollary is not asserted.  When
$p(S,T)=d$, the bridge is equality; when $d=0$, the exact boundary conclusion
of Lemma~\ref{lem:step-001-zero} remains the baseline rather than being
replaced by a positive-dimensional surrogate.
\end{proof}

\begin{proof}[Derivation of the dimension bound]
Proposition~\ref{prop:step-005-routing} first handles every boundary branch.
If $\mathcal H=\varnothing$, the exact convention gives
$\operatorname{dc}(\mathcal H)=0$, and no logarithmic budget or feature tuple
is needed.  If $\mathcal H\neq\varnothing$ and $d=0$,
Lemma~\ref{lem:step-001-zero} gives
$\operatorname{dc}(\mathcal H)=0=7TSd$ with the original tie label.

On the remaining branch, Proposition~\ref{prop:step-001-architecture}
supplies $d\geq1$ and $T,S\geq1$.
Proposition~\ref{prop:step-003-budget} supplies the exact integer
$r=\lceil\log_2(2|\mathcal H|)\rceil\leq7TS$, while
Proposition~\ref{prop:step-004-covering} supplies one deterministic tuple that
works for every target.  Lemma~\ref{lem:step-005-score} concatenates that
tuple into one map $\Phi$ chosen before any target-specific separator.  Its
pointwise score identity is exact, so the fixed $\operatorname{sgn}_{\tau}$,
including zero-score ties, returns each target exactly.
Proposition~\ref{prop:step-005-dimension} then invokes the definition of
$\operatorname{dc}$ and the explicit comparison $rd\leq7TSd$.  Thus every
setup under the three assumptions obeys
\[
\operatorname{dc}(\mathcal H)\leq7TSd.
\]
The common map is target-independent: the tuple is fixed once, whereas only
$i(h)$, $w_h$, and $u_h$ depend on $h$, exactly as permitted in the
definition of $\operatorname{dc}$.  Corollary
\ref{cor:step-005-polynomial-bridge} performs the optional specialization only
when a separate explicit inequality $d\leq p(S,T)$ is available.
When $r=1$, the direct sum is the successful block itself.  When
$p(S,T)=d$, the polynomial comparison is equality, and when $d=0$ the exact
zero-dimensional conclusion is retained instead of replaced by a positive
surrogate.
\end{proof}

\subsection{Proof of the Main Theorem}
\label{app:main-proof}

\begin{proof}[Proof of Theorem~\ref{thm:main}]
Under Assumptions~\ref{assump:source-regime},
\ref{assump:universal-sgd-success}, and
\ref{assump:tie-resolved-confident-map},
Proposition~\ref{prop:step-005-routing} gives an exhaustive case split.  Its
first two alternatives establish the result directly when
$\mathcal H=\varnothing$ or $d=0$.  In its third alternative,
Proposition~\ref{prop:step-005-dimension}, whose inputs are supplied by
Propositions~\ref{prop:step-002-vc}, \ref{prop:step-003-budget}, and
\ref{prop:step-004-covering}, gives
\[
\operatorname{dc}(\mathcal H)\leq rd\leq7TSd.
\]
Every branch therefore satisfies the claimed deterministic bound.  The only
use of Assumption~\ref{assump:universal-sgd-success} is the VC ceiling in
Proposition~\ref{prop:step-002-vc}; the target-independent confident-map
premise remains a separate explicit condition and is not derived from SGD.
\end{proof}
